% Methods section draft for PlanetProfile-genai manuscript
% Generated by manuscript-writer agent
% Date: 2026-04-23
%
% INSTRUCTIONS: Copy this content into main.tex Section 2 (Methods)
% Add references to sample.bib as needed

\section{Methods}

\subsection{PlanetProfile Framework}

We use PlanetProfile \citep{vance2018planetprofile}, an open-source Python framework for constructing one-dimensional interior structure models of ocean worlds and icy satellites. PlanetProfile employs self-consistent thermodynamics for fluid, rock, and mineral phases, propagating planetary properties layer-by-layer from the surface to the core. The framework incorporates:

\begin{itemize}
    \item Ocean equation of state calculations using SeaFreeze \citep{journaux2020seafreeze} for water and aqueous solutions (NaCl, MgSO4, seawater, or custom compositions via Reaktoro \citep{leal2017reaktoro})
    \item High-pressure ice thermodynamics based on \citet{journaux2020gibbs}, including phases Ih, II, III, V, VI, and clathrate hydrates
    \item Silicate and metallic core properties from PerpleX \citep{connolly2005perplex} equation of state tables
    \item Forward modeling of tidal Love numbers via PyALMA3 \citep{beuthe2013pyalma3} or TidalPy \citep{tobie2024tidalpy}
    \item Magnetic induction responses using MoonMag \citep{styczinski2021moonmag}
\end{itemize}

For each model body, we specify the bulk planetary mass and radius, surface temperature and gravity, and core composition. PlanetProfile then iteratively solves for the pressure, temperature, and thermodynamic properties at each depth, ensuring hydrostatic equilibrium and energy balance throughout the interior.

\subsection{Two-Phase Convection in High-Pressure Ice}

\subsubsection{Kalousova \& Sotin (2018) Parameterization}

We implement the two-phase convection model of \citet{kalousova2018melting} for high-pressure ice layers (phases III, V, and VI). This parameterization extends beyond traditional solid-state convection (e.g., \citealt{deschamps2001thermal} for ice Ih) by accounting for partial melting and liquid water transport within convecting HP ice.

The model characterizes four exchange regimes between the overlying ocean and underlying silicate mantle, determined by a modified Rayleigh number:

\begin{equation}
\text{Ra}^* = \frac{\alpha_0 \rho_0 g^2 c_p \Delta T_c H_c^3}{\kappa \nu_0}
\label{eq:rastar}
\end{equation}

\noindent where $\alpha_0$ is the thermal expansivity, $\rho_0$ the density, $g$ the gravitational acceleration, $c_p$ the heat capacity, $\Delta T_c = T_{\text{melt}}(\text{base}) - T_{\text{interior}}$ the temperature contrast across the convecting layer, $H_c$ the layer thickness, $\kappa$ the thermal diffusivity, and $\nu_0$ the reference kinematic viscosity at the melting temperature.

The critical Rayleigh number for temperate layer formation is given by:

\begin{equation}
\text{Ra}^*_c = 19.965 \times 10^3 \times (q_s)^{3.690}
\label{eq:racrit}
\end{equation}

\noindent where $q_s$ is the heat flux from the silicate mantle in mW~m$^{-2}$.

When $\text{Ra}^* > \text{Ra}^*_c$, a temperate (partially molten) layer forms at the ocean-ice interface with thickness:

\begin{equation}
H_t[\text{km}] = (0.145 \times 10^{-3} \times q_s[\text{mW m}^{-2}] + 0.015) \times \mu_0[\text{Pa s}]^{0.21}
\label{eq:ht}
\end{equation}

\noindent where $\mu_0 = \rho_0 \nu_0$ is the reference dynamic viscosity.

The thermal boundary layer (TBL) at the ice-silicate interface scales as:

\begin{equation}
\delta = 2.746 \times (\text{Ra}^*)^{-0.271} \times H_c
\label{eq:delta}
\end{equation}

\subsubsection{Implementation Details}

Our implementation in PlanetProfile distinguishes between ice Ih (which uses the \citealt{deschamps2001thermal} solid-state convection model) and high-pressure ices III, V, and VI (which use the \citealt{kalousova2018melting} two-phase model when enabled via the \texttt{KALOUSOVA\_CONVECTION} flag).

Key aspects of the implementation include:

\begin{enumerate}
    \item \textbf{Temperature profile}: The convecting interior follows the pressure-dependent melting curve for the given ice phase, consistent with vigorous two-phase convection. This differs from the well-mixed interior temperature assumption in solid-state convection models.

    \item \textbf{Reference viscosity}: We use phase-specific reference viscosities $\eta_{\text{melt}}$ evaluated at the melting temperature for each HP ice phase, derived from laboratory measurements compiled in \citet{durham2010ice}.

    \item \textbf{Heat flux}: The model accounts for heat flux from the silicate mantle, with optional contributions from tidal heating within the HP ice layer (though the latter is not included in the baseline \citealt{kalousova2018melting} formulation).

    \item \textbf{Triple-point boundaries}: Phase transitions between ice III, V, and VI are determined using triple-point temperatures and pressures from \citet{journaux2020gibbs}: $T_{\text{III-V-L}} = 254.0$~K at $P = 350.0$~MPa, and $T_{\text{V-VI-L}} = 272.0$~K at $P = 632.0$~MPa.

    \item \textbf{Melt fraction}: When a temperate layer is detected ($\text{Ra}^* > \text{Ra}^*_c$), we store a nominal melt fraction of 0.01. The \citet{kalousova2018melting} parameterization does not provide an explicit melt fraction calculation; this value serves as a placeholder for future refinement with a full two-phase flow model.
\end{enumerate}

The convection state for each HP ice phase is stored in the planetary structure, including the convecting interior temperature $T_{\text{conv}}$, temperate layer thickness $H_t$ (if present), TBL thickness $\delta$, and the Rayleigh numbers $\text{Ra}^*$ and $\text{Ra}^*_c$.

\subsection{Clathrate Insulation Model}

% TODO: Add clathrate insulation description
% Placeholder: Describe clathrate hydrate model in ice Ih shell
% Include thermal conductivity reduction, stability field
% Reference: Kalousova & Sotin (2020) Titan paper?

\subsection{Rheology and Viscosity Models}

\subsubsection{High-Pressure Ice Viscosity}

% TODO: Add composite rheology description
% Placeholder: Diffusion creep + dislocation creep
% Arrhenius temperature dependence
% Grain size effects
% Reference: Durham et al. (2010)?

\subsubsection{Silicate Viscosity}

% TODO: Add silicate viscosity model
% Placeholder: Temperature and pressure-dependent viscosity
% Activation energy parameters
% Critical for heat flux and ocean stability

\subsection{Titan Model Configuration}

% TODO: Add Titan-specific configuration
% Placeholder:
% - Bulk properties (mass, radius, gravity)
% - Ocean composition (NH3-H2O?)
% - Clathrate lid thickness
% - With/without ocean scenarios
% - Heat source variations
% Reference specific PPTest files or configurations

% NOTES FOR TEAM:
% - Need input on clathrate model details
% - Need viscosity parameter values to include
% - Need specific Titan configuration details
% - Need to verify which PPTest files correspond to Titan scenarios
% - Total word count so far: ~800 words
% - Target: ~1500-2000 words for full Methods section
